whence the result by (6.1.2), after passage to the quotient.

If @@M0@@, @@M1@@, so @@M2@@ is a homomorphism of @@M3@@-algebras; moreover, it is compatible with the augmentations.

Finally, it remains to see that it is a @@M4@@-homomorphism. We must therefore prove the commutativity of the diagram:
@@M5@@

By additivity, it suffices to prove this for the elements of @@M6@@, for which it is evident, and for the elements of zero augmentation of @@M7@@, for which it follows, by the definition of @@M8@@ and (6.1.4), from that of diagram (3.7.2).

\sgaStableTarget{sga6.e05.prop.6-9}\textbf{Proposition 6.9.} \uline{Let} @@M9@@ \uline{be an augmented} @@M10@@-@@M11@@-\uline{algebra}, @@M12@@; \uline{then}:
@@M13@@
\uline{i}.\uline{e}.
@@M14@@

We must see that these two elements give the same image in @@M15@@, that is, that the @@M16@@-th Chern class of @@M17@@ is the image of @@M18@@.

We can write @@M19@@ as a sum @@M20@@, with @@M21@@, @@M22@@ for every @@M23@@, and @@M24@@. We then have:
@@M25@@
Since @@M26@@ is a @@M27@@-homomorphism, we can write:
@@M28@@
